\part{What an operating system is}

\chapter{The minimum pieces}

\begin{goals}
  \item Which problems an operating system solves, and which it does not
  \item The seven pieces without which there is no operating system
  \item In what order they must be built, and why that order is mandatory
\end{goals}

\section{The problem it solves}

Imagine you want to run two programs at once on the computer from Chapter 1. With no operating
system, you have to answer a pile of awkward questions yourself:

\begin{itemize}
  \item Where does each program load? If both expect to sit at \hexa{400000}, they clobber each
        other.
  \item How do they share the processor? There is only one \reg{EIP}.
  \item If one enters an infinite loop, how do I get control back?
  \item If one writes to a wrong address, how do I stop it wrecking the other?
  \item Both want to write to the screen. Who wins?
  \item Both want to save a file. How do they know which disk sectors are free?
\end{itemize}

An operating system is exactly the answer to that list. It is not ``the program that starts when
you power on''; it is \hi{the referee that shares limited hardware between programs that do
not know each other and do not trust each other}.

\begin{concept}{The three jobs of a kernel}
Everything a kernel does falls into one of three categories:
\begin{enumerate}
  \item \hi{Abstract}: turn ``disk sectors'' into ``files'', ``physical pages'' into ``your
        own memory'', ``one processor'' into ``every program has its own''.
  \item \hi{Share}: decide who uses what, and for how long.
  \item \hi{Protect}: guarantee that one program cannot damage another or the kernel itself.
\end{enumerate}
Every technical chapter in this course implements one of the three, for one specific resource.
\end{concept}

\section{The seven pieces}

A minimal but genuine operating system --- able to run several isolated programs --- needs these
seven pieces. Not one fewer.

\begin{center}
\begin{tikzpicture}[font=\sffamily\scriptsize,
  p/.style={draw=accent,fill=accentlight,rounded corners=3pt,minimum width=3.9cm,
            minimum height=1.15cm,align=center,text=ink}]

  \node[p] (a) {\textbf{1. Boot}\\{\tiny take control of the hardware}};
  \node[p,below=0.45cm of a] (b) {\textbf{2. Protected mode}\\{\tiny enable privileges and 32 bits}};
  \node[p,below=0.45cm of b] (c) {\textbf{3. Output}\\{\tiny be able to see what happens}};
  \node[p,below=0.45cm of c] (d) {\textbf{4. Interrupts}\\{\tiny react to the outside world}};

  \node[p,right=1.6cm of a] (e) {\textbf{5. Memory management}\\{\tiny share and isolate RAM}};
  \node[p,below=0.45cm of e] (f) {\textbf{6. Processes}\\{\tiny several programs at once}};
  \node[p,below=0.45cm of f] (g) {\textbf{7. Storage}\\{\tiny files that survive}};
  \node[draw=green,fill=greenlight,rounded corners=3pt,minimum width=3.9cm,
        minimum height=1.15cm,align=center,below=0.45cm of g] (h)
        {\textbf{\color{green}User space}\\{\tiny unprivileged programs}};

  \draw[-{Stealth[length=2.5mm]},accent,thick] (a)--(b);
  \draw[-{Stealth[length=2.5mm]},accent,thick] (b)--(c);
  \draw[-{Stealth[length=2.5mm]},accent,thick] (c)--(d);
  \draw[-{Stealth[length=2.5mm]},accent,thick] (d.east) -- ++(0.4,0) |- (e.west);
  \draw[-{Stealth[length=2.5mm]},accent,thick] (e)--(f);
  \draw[-{Stealth[length=2.5mm]},accent,thick] (f)--(g);
  \draw[-{Stealth[length=2.5mm]},green,thick] (g)--(h);
\end{tikzpicture}
\end{center}

Let us walk through them, without yet worrying about implementation.

\subsection{1. Boot: taking control}

Somebody has to go first. The firmware loads 512 bytes from disk and jumps into them; those 512
bytes have to load the rest. It is a problem of pulling yourself up by your own bootstraps ---
hence \emph{bootstrapping}, and hence \emph{boot}.

\subsection{2. Protected mode: unlocking the processor}

The processor starts up pretending to be a chip from 1978, for compatibility. It sees only 1 MB of
memory, works in 16 bits and has no protection at all. It must be switched explicitly. Without
this, pieces 5 to 7 are literally impossible.

\subsection{3. Output: being able to see what happens}

This looks trivial and is not. Debugging an operating system with no way to print is like fixing
an engine with your eyes shut. It is the first thing built after you have a working execution
environment, and the reason is purely practical.

\subsection{4. Interrupts: reacting}

Without interrupts, a program that enters an infinite loop keeps the processor forever. Interrupts
are the mechanism by which hardware can \emph{take} control away from running code and hand it to
the kernel. They are the basis of:

\begin{itemize}
  \item multitasking (the timer interrupts and the kernel switches process),
  \item devices (the keyboard interrupts when a key is pressed),
  \item error handling (dividing by zero interrupts),
  \item system calls (a program interrupts itself to ask for something).
\end{itemize}

\begin{concept}{Interrupts are the axis}
If you had to keep one mechanism to understand an operating system, this would be it. Almost
everything interesting a kernel does happens \emph{inside an interrupt handler}. In lytlnyblOS,
process switching, keyboard input and every system call all go through one.
\end{concept}

\subsection{5. Memory management: sharing and isolating}

It comes in three stacked layers:

\begin{center}
\begin{tabularx}{\textwidth}{@{}l l X@{}}
\toprule
\textbf{Layer} & \textbf{Unit} & \textbf{Answers the question} \\
\midrule
Physical memory & 4 KB page & Which chunks of RAM are free? \\
Virtual memory  & 4 KB page & What does each program see when it looks at an address? \\
Heap            & bytes     & How do I ask for 17 bytes without wasting a whole page? \\
\bottomrule
\end{tabularx}
\end{center}

The middle layer, virtual memory, is the hardest and the most important: it is what makes it
possible for two programs to use the same address \hexa{400000} without colliding. It gets a
chapter of its own, explained very slowly.

\subsection{6. Processes: several programs at once}

A \term{process} is a program \emph{in execution} plus all its state: its register values, its
stack, its memory, its open files. The kernel stores that state in a structure and swaps it dozens
of times a second, creating the illusion of simultaneity.

\subsection{7. Storage: making things survive}

RAM is erased when you power off. For a file to exist tomorrow, it has to be written to disk; and
for it to be \emph{findable} tomorrow, the disk needs a structure imposed on it: a
\term{filesystem}.

\section{Why this order and no other}

The order is not pedagogical: it is a chain of real technical dependencies.

\begin{center}
\begin{tabularx}{\textwidth}{@{}l X@{}}
\toprule
\textbf{To build\ldots} & \textbf{you first need\ldots} \\
\midrule
Protected mode       & to have booted and have the code in memory \\
VGA text mode        & protected mode (to address above 1 MB comfortably) \\
Interrupts           & protected mode (the interrupt table is different) and output (to see faults) \\
Timer                & working interrupts and a remapped PIC \\
Physical memory      & the memory map the bootloader stored \\
Virtual memory       & the physical manager (to request pages for the tables) \\
Kernel heap          & virtual memory (to map more pages as it grows) \\
Processes            & the heap (to allocate structures) and interrupts (to switch) \\
User space           & processes, virtual memory and a TSS \\
System calls         & user space and interrupts \\
Filesystem           & a disk driver \\
The shell            & absolutely everything above \\
\bottomrule
\end{tabularx}
\end{center}

\begin{warn}{The cost of skipping a link}
It is tempting to try loading user programs before paging is properly sorted, or to write the
filesystem before the heap. You can, but you pay for it: the code fills up with special cases that
later have to be unpicked. This course follows the order the project is actually built in, and
that is not a coincidence.
\end{warn}

% ============================================================
\chapter{Kernel and user space}

\begin{goals}
  \item What ``user space'' means exactly
  \item What a system call is and why it is not an ordinary function call
  \item What has to happen to cross the boundary in each direction
\end{goals}

\section{Two worlds}

We saw that the processor has privilege rings. An operating system uses them to split software into
two worlds with different rules:

\begin{center}
\begin{tikzpicture}[font=\sffamily\scriptsize]
  \fill[greenlight,draw=green,rounded corners=4pt] (0,2.3) rectangle (11,4.6);
  \node[text=green,font=\sffamily\bfseries] at (1.75,4.25) {USER SPACE (ring 3)};
  \node[draw=green,fill=white,rounded corners=2pt,minimum width=2.2cm,minimum height=0.85cm] at (1.8,3.1) {shell};
  \node[draw=green,fill=white,rounded corners=2pt,minimum width=2.2cm,minimum height=0.85cm] at (4.6,3.1) {user\_test};
  \node[draw=green,fill=white,rounded corners=2pt,minimum width=2.2cm,minimum height=0.85cm] at (7.4,3.1) {your program};
  \node[text=gray,font=\sffamily\tiny,align=left] at (9.8,3.1) {each with its\\own private\\memory};

  \fill[redlight,draw=red,rounded corners=4pt] (0,0) rectangle (11,1.75);
  \node[text=red,font=\sffamily\bfseries] at (1.3,1.42) {KERNEL (ring 0)};
  \node[draw=red,fill=white,rounded corners=2pt,minimum width=1.9cm,minimum height=0.7cm] at (1.6,0.62) {memory};
  \node[draw=red,fill=white,rounded corners=2pt,minimum width=1.9cm,minimum height=0.7cm] at (3.9,0.62) {processes};
  \node[draw=red,fill=white,rounded corners=2pt,minimum width=1.9cm,minimum height=0.7cm] at (6.2,0.62) {files};
  \node[draw=red,fill=white,rounded corners=2pt,minimum width=1.9cm,minimum height=0.7cm] at (8.5,0.62) {drivers};

  \draw[grayborder,line width=1.5pt,dashed] (0,2.02) -- (11,2.02);
  \node[fill=white,text=accent,font=\sffamily\bfseries\scriptsize] at (5.5,2.02) {int 0x80 --- the only door};
\end{tikzpicture}
\end{center}

The dashed line is the boundary. It is a \emph{real} boundary, enforced by the processor, not a
convention. A user program cannot call a kernel function even if it knows the address: the page is
marked supervisor-only and the attempt raises an exception.

\section{The only door}

If a program cannot call the kernel, how does it ask it to print something? Through the one
mechanism the processor allows: \hi{raising an interrupt on purpose}.

\begin{center}
\begin{tikzpicture}[font=\sffamily\scriptsize,
  p/.style={draw=grayborder,fill=graylight,rounded corners=2pt,minimum width=5.4cm,minimum height=0.72cm,align=left,inner xsep=8pt}]
  \node[p,fill=greenlight,draw=green] (a) {1. The program puts the service number in EAX};
  \node[p,fill=greenlight,draw=green,below=0.28cm of a] (b) {2. It puts arguments in EBX, ECX, EDX};
  \node[p,fill=greenlight,draw=green,below=0.28cm of b] (c) {3. It executes \texttt{int 0x80}};
  \node[p,fill=amberlight,draw=amber,below=0.28cm of c] (d) {4. \textbf{The processor switches to ring 0}};
  \node[p,fill=redlight,draw=red,below=0.28cm of d] (e) {5. It jumps where the IDT says};
  \node[p,fill=redlight,draw=red,below=0.28cm of e] (f) {6. The kernel reads EAX and does the work};
  \node[p,fill=redlight,draw=red,below=0.28cm of f] (g) {7. It leaves the result in EAX};
  \node[p,fill=amberlight,draw=amber,below=0.28cm of g] (h) {8. \texttt{iret}: \textbf{back to ring 3}};
  \foreach \x/\y in {a/b,b/c,c/d,d/e,e/f,f/g,g/h} \draw[-{Stealth[length=2mm]},gray] (\x)--(\y);
\end{tikzpicture}
\end{center}

Steps 4 and 8 are the important ones, and the \emph{hardware} performs them, not software. It is
the processor that decides to raise privilege, and it only does so by jumping to an address the
kernel registered in advance. The program does not choose where it lands: it chooses \emph{what it
asks for}, not \emph{where it goes}.

\begin{concept}{Why this is safe}
A malicious program cannot ``jump into the kernel at a convenient line'' because the only entry is
the one the interrupt table defines. The kernel controls the landing point. All the attacker
controls are the numbers in the registers --- which is why \hi{validating those numbers is
the single most important thing a system call handler does}. Part XIII shows what happens when
that validation is missing.
\end{concept}

\section{The ten services of lytlnyblOS}

Our system offers ten calls. That is a ridiculously small number (Linux has over three hundred)
but enough for a complete shell.

\begin{center}
\begin{tabularx}{\textwidth}{@{}l l l X@{}}
\toprule
\textbf{EAX} & \textbf{Name} & \textbf{Arguments} & \textbf{What it does} \\
\midrule
\hexa{01} & \cod{exit}    & ---               & Terminates the process and wakes its parent \\
\hexa{02} & \cod{getpid}  & ---               & Returns the process identifier \\
\hexa{03} & \cod{yield}   & ---               & Voluntarily gives up the processor \\
\hexa{04} & \cod{sleep}   & ticks             & Sleeps for $n$ timer ticks \\
\hexa{05} & \cod{write}   & fd, buffer, n     & Writes to the screen or to a file \\
\hexa{06} & \cod{read}    & fd, buffer, n     & Reads from the keyboard or from a file \\
\hexa{07} & \cod{sbrk}    & increment         & Grows the process heap \\
\hexa{08} & \cod{fsops}   & operation, path   & \cod{ls}, \cod{mkdir}, \cod{touch}, \cod{rm}, \cod{run} \\
\hexa{09} & \cod{open}    & path              & Opens a file and returns a descriptor \\
\hexa{0A} & \cod{close}   & fd                & Closes a descriptor \\
\bottomrule
\end{tabularx}
\end{center}

Notice that \cod{write} serves both the screen and files. That is an idea inherited from Unix:
\hi{everything is a descriptor}. The program does not know whether 1 is the screen or a
file; it only knows it is writing to ``descriptor 1''. The kernel decides what that means.

\exercise{Look at the table and work out the minimum set of calls the shell's \cod{cat file.txt}
command needs, and in what order. (The answer is in Part XII, but try first.)}

% ============================================================
\chapter{Map of the whole system}

\begin{goals}
  \item See the complete boot journey end to end
  \item Know which project file implements each stage
  \item Have a mental map to come back to whenever you get lost
\end{goals}

This chapter is the reference to return to whenever you wonder ``where in the system am I?''. You
do not need to understand it all now; you need to know it exists.

\section{The journey, end to end}

\begin{center}
\begin{tikzpicture}[font=\sffamily\scriptsize,
  e/.style={draw=accent,fill=accentlight,rounded corners=3pt,minimum width=5.0cm,
            minimum height=0.85cm,align=center},
  f/.style={font=\ttfamily\tiny,text=purple,anchor=west}]

  \node[e,fill=graylight,draw=gray] (bios) {Firmware (BIOS): checks the hardware};
  \node[f] at (5.2,0) {(in ROM, not ours)};

  \node[e,below=0.38cm of bios] (boot) {Boot sector: loads the kernel};
  \node[f] at (5.2,-1.23) {boot/bootloader.asm};

  \node[e,below=0.38cm of boot] (intro) {Entry: GDT and protected mode};
  \node[f] at (5.2,-2.46) {kernel/kernel\_intro.asm};

  \node[e,below=0.38cm of intro] (kmain) {\textbf{kernel\_main(): brings up subsystems}};
  \node[f] at (5.2,-3.69) {kernel/kernel\_main.c};

  \node[e,below=0.38cm of kmain,fill=amberlight,draw=amber] (idt) {idt\_init(): interrupts and PIC};
  \node[f] at (5.2,-4.92) {kernel/interrupts.c};

  \node[e,below=0.38cm of idt,fill=amberlight,draw=amber] (mem) {init\_pmm/vmm/heap: memory};
  \node[f] at (5.2,-6.15) {memory/*.c};

  \node[e,below=0.38cm of mem,fill=amberlight,draw=amber] (proc) {init\_procman + init\_tss: processes};
  \node[f] at (5.2,-7.38) {tasks/procman.c};

  \node[e,below=0.38cm of proc,fill=amberlight,draw=amber] (drv) {timer\_init + keyboard\_init + fs};
  \node[f] at (5.2,-8.61) {kernel/drivers/, filesystem/};

  \node[e,below=0.38cm of drv,fill=greenlight,draw=green] (load) {load\_program("/bin/shell")};
  \node[f] at (5.2,-9.84) {tasks/loader.c};

  \node[e,below=0.38cm of load,fill=greenlight,draw=green] (sh) {\textbf{The shell runs in ring 3}};
  \node[f] at (5.2,-11.07) {user/programs/shell.c};

  \foreach \x/\y in {bios/boot,boot/intro,intro/kmain,kmain/idt,idt/mem,mem/proc,proc/drv,drv/load,load/sh}
    \draw[-{Stealth[length=2.2mm]},accent,thick] (\x)--(\y);
\end{tikzpicture}
\end{center}

\section{What each directory holds}

\begin{center}
\begin{tabularx}{\textwidth}{@{}l X@{}}
\toprule
\textbf{Directory} & \textbf{Contents} \\
\midrule
\filepath{boot/}        & The boot sector. 512 bytes of 16-bit assembly. \\
\filepath{kernel/}      & Protected-mode entry, interrupts, utilities and \cod{kernel\_main}. \\
\filepath{kernel/drivers/} & VGA text mode, PIT timer, PS/2 keyboard and ATA disk. \\
\filepath{memory/}      & Physical manager (bitmap), virtual manager (paging) and heap. \\
\filepath{tasks/}       & Processes, context, scheduler and program loader. \\
\filepath{filesystem/}  & The filesystem in two layers: \cod{fs.c} (low level) and \cod{fs\_manager.c} (descriptors and paths). \\
\filepath{user/libc/}   & The minimal standard library linked into every user program. \\
\filepath{user/programs/} & The programs: the shell and a test program. \\
\filepath{tools/}       & \cod{mkfs.c}: compiled for \emph{your} machine, not for the system. \\
\filepath{include/}     & All headers, mirroring the structure above. \\
\filepath{rootfs/}      & The files \cod{mkfs} injects into the disk image. \\
\bottomrule
\end{tabularx}
\end{center}

\section{What the build produces}

\cod{make} generates a single 10 MB disk image laid out like this:

\begin{center}
\begin{tikzpicture}[font=\sffamily\scriptsize,x=1cm,y=1cm]
  \def\w{11}
  \fill[purplelight,draw=purple] (0,0) rectangle (1.3,1.1);
  \node[align=center,font=\tiny] at (0.65,0.55) {boot\\512 B};
  \node[font=\ttfamily\tiny,text=gray] at (0.65,-0.32) {sector 0};

  \fill[accentlight,draw=accent] (1.3,0) rectangle (5.2,1.1);
  \node[align=center,font=\tiny] at (3.25,0.55) {\textbf{kernel.bin}\\$\sim$23 KB (sectors 1--45)};
  \node[font=\ttfamily\tiny,text=gray] at (1.4,-0.32) {sector 1};

  \fill[graylight,draw=grayborder] (5.2,0) rectangle (6.4,1.1);
  \node[align=center,font=\tiny] at (5.8,0.55) {free};

  \fill[greenlight,draw=green] (6.4,0) rectangle (\w,1.1);
  \node[align=center,font=\tiny] at ({(6.4+\w)/2},0.55) {\textbf{filesystem}\\superblock, inodes, data};
  \node[font=\ttfamily\tiny,text=gray] at (6.5,-0.32) {sector 70};
\end{tikzpicture}
\end{center}

\begin{warn}{The 60-sector ceiling}
The boot sector asks the BIOS for exactly 60 sectors (30,720 bytes). The kernel currently occupies
around 23 KB, so there is headroom --- but not unlimited. If the kernel exceeded that size, the
BIOS would load only part of it and the system would fail in a baffling way: no error message,
just executing garbage. It is a limit worth remembering before adding large code to the kernel.
\end{warn}

\begin{summary}
You now know which pieces exist, what each is for, and where its code lives. From here we build, in
order, starting at the literal beginning: the first 512 bytes.
\end{summary}
